\section{SayCan: el LLM planifica y la affordance decide si se puede hacer}

RT-1 amplió la skill library, pero no resolvió cómo una instrucción como «limpia la mesa» se descompone en varias habilidades sucesivas. SayCan usa el modelo de lenguaje como high-level planner y, al mismo tiempo, emplea el value / affordance model del robot para restringir si cada paso es ejecutable en el estado actual; así conecta el Internet-scale common sense con el grounded control.

\subsection{De Skill Learning a Planning}

El tercer punto de la Agenda pasa a SayCan. A continuación, la timeline de la clase vincula las multi-task robot policies de 2021--2022 con el LLM-powered planning de 2022--2023: las acciones de bajo nivel siguen proviniendo de robot data, mientras que la composición de alto nivel empieza a apoyarse en modelos de lenguaje.

\lecturefigure{slide-23-agenda-saycan.jpg}{Agenda: SayCan entra en la capa de language-conditioned planning}{Diapositiva V023 recuperada de la grabación oficial; Stanford CS25 Lecture 15}

\lecturefigure{slide-24-robotics-foundation-timeline.jpg}{Timeline de la investigación en robótica: de escalar operaciones reales a acelerar los sistemas robóticos con LLM}{Diapositiva V024 recuperada de la grabación oficial; Stanford CS25 Lecture 15}

\subsection{Dos problemas: Fixed Commands y No Grounding}

Primero, el robot solo puede invocar un conjunto limitado de habilidades, por ejemplo `find apple`, `pick apple`, `go table`; el LLM tiene que aprender a «hablar el lenguaje del robot». Segundo, el LLM no tiene cámara ni dinámica: no sabe si la manzana existe ni si la pinza puede sujetarla; el sistema debe añadir physical affordance grounding al conocimiento previo del lenguaje.

\lecturefigure{slide-25-language-models-robotics-challenges.jpg}{Language models for robotics: skill vocabulary fijo y falta de real-world grounding}{Diapositiva V025 recuperada de la grabación oficial; Stanford CS25 Lecture 15}

\begin{warningbox}{El planner no puede crear Skills que no existen}
El LLM puede descomponer un objetivo de manera impecable, pero si la low-level library no incluye «abrir el cajón» o «limpiar un líquido», el planner solo puede producir texto no ejecutable. El límite superior del sistema lo determina la cobertura de las grounded skills; este es también el SayCan bottleneck que el ponente señaló explícitamente en la sesión de preguntas y respuestas.
\end{warningbox}

\subsection{SayCan Score: el producto de Useful y Possible}

Para una habilidad candidata $k$, un objetivo de alto nivel $g$ y el estado actual $s$, SayCan puede escribirse como
\begin{equation}
\operatorname{Score}(k\mid g,s)
\propto p_{\mathrm{LM}}(k\mid g,h)
\cdot p_{\mathrm{aff}}(\mathrm{success}\mid s,k),
\end{equation}
donde $h$ son los pasos ya ejecutados. El primer término pregunta «¿este paso sirve para cumplir el objetivo expresado en lenguaje?», y el segundo, «¿el robot puede hacerlo ahora?». En log space ambos términos se suman, lo que evita que una acción plausible en lenguaje pero físicamente imposible quede en primer lugar.

\lecturefigure{slide-26-saycan-architecture.jpg}{SayCan: la usefulness del LLM y la robotic affordance eligen conjuntamente la siguiente habilidad}{Diapositiva V026 recuperada de la grabación oficial; Ahn et al., 2022}

\begin{knowledgebox}{Lectura de la figura: ordenamiento de candidatos, no control generado libremente}
La tabla central enumera habilidades como `find apple` y `pick apple`. La puntuación del LLM favorece los pasos coherentes con el objetivo, la puntuación de affordance filtra los que no son ejecutables en ese momento, y el sistema elige el de mayor combined score; tras ejecutarlo, vuelve a ordenar. El controller de bajo nivel no lo acciona directamente ningún token del LLM.
\end{knowledgebox}

\subsection{Experimentos: Planning y Execution se puntúan por separado}

El diagrama de arquitectura muestra cómo SayCan combina usefulness y affordance, pero aún hace falta distinguir dos afirmaciones diferentes: «el plan es razonable en lenguaje» y «el robot completa finalmente la tarea». Esta sección usa métricas independientes para examinar la contribución del ordenamiento de alto nivel y la de la ejecución de bajo nivel; al leer los resultados, conviene comparar primero la diferencia entre planning y execution, y después revisar la ablation que retira el grounding. Los experimentos demuestran que la puntuación de ejecutabilidad en el entorno es indispensable, pero la tasa de éxito global por sí sola no permite determinar si los fallos restantes provienen de la percepción, del control o de la skill coverage.

En la clase se reportan 101 long-horizon instructions, con una tasa de éxito de planificación de aproximadamente 84\% y una tasa de éxito de ejecución de aproximadamente 74\%, y se requiere encadenar más de diez navigation / manipulation skills. Al retirar el grounding, el desempeño cae casi a la mitad, lo que indica que la plausibility lingüística no sustituye la ejecutabilidad en el entorno.

\lecturefigure{slide-27-saycan-experiment-results.jpg}{SayCan experiment overview: 84\% de planning, 74\% de execution y grounding ablation}{Diapositiva V027 recuperada de la grabación oficial; Ahn et al., 2022}

\begin{warningbox}{Planning Success no equivale a Task Success}
Una secuencia de pasos puede ser semánticamente correcta y aun así no ejecutarse por fallos de percepción, de sujeción o de navegación. Reportar planning y execution por separado permite localizar el cuello de botella; si solo se reporta la tasa final de la tarea, no se puede saber si el fallo proviene del LLM, del affordance model o de la low-level skill.
\end{warningbox}

\subsection{PaLM-SayCan: cómo se transmite el progreso de un modelo externo}

En la slide 28, el eje horizontal es el language model size y el eje vertical, la planning performance. PaLM-SayCan mejora conforme aumenta la escala del modelo, lo que indica que, con los datos del robot sin cambios, el progreso de un LLM externo puede transmitirse al sistema a través de la interfaz del planner. El punto azul de FLAN-SayCan recuerda también que el instruction tuning y la familia del modelo se mezclan en la comparación de scale.

\lecturefigure{slide-28-palm-saycan-scaling.jpg}{PaLM-SayCan: un LLM más capaz mejora el desempeño de la planificación de alto nivel}{Diapositiva V028 recuperada de la grabación oficial; Stanford CS25 Lecture 15}

Un mejor prompt y el chain-of-thought permiten además atender solicitudes como «busca una fruta que no sea manzana» o peticiones de limpieza de varios pasos. Los ejemplos de la derecha convierten la explicación en lenguaje natural en una secuencia de habilidades invocables, lo que muestra la decomposition semántica del LLM; sin embargo, cada secuencia todavía debe pasar por la affordance score antes de llegar al robot.

\lecturefigure{slide-29-palm-saycan-prompting.jpg}{PaLM-SayCan prompting: CoT e indicaciones más elaboradas generan secuencias de habilidades de varios pasos}{Diapositiva V029 recuperada de la grabación oficial; Stanford CS25 Lecture 15}

\teachervoice{Lo que el ponente más valora de SayCan es su «leverage»: sin volver a recolectar trayectorias del robot, basta con sustituir el LLM o el prompt por uno mejor para que la capacidad de planificación aumente. Pero también subraya que este leverage solo actúa sobre la skill selection y no puede suplir acciones que el low-level controller simplemente no sabe realizar.}

\subsection{Resumen de la sección}

SayCan conecta el conocimiento previo del lenguaje con la ejecución física mediante el producto de la LM usefulness y la affordance probability. Demuestra que el progreso de un LLM externo puede transmitirse a la planning del robot y, a la vez, deja muy claro el límite superior del sistema: la grounded skill library, el affordance model y la ejecución de bajo nivel siguen siendo cuellos de botella que no se pueden omitir.
